% BEGIN CR28_CONSTITUTIVE_INVERSE
\section{Functional recovery, composition and the Catalan moment}
\label{cr28:sec:restitucion-constitutiva}

The incidence response \(90/120\) preserves more information than an isolated scalar value. In the same notation as \eqref{eq:three-four-completion}, write
\[
 f(s)=\frac{1+s+s^2}{1+s+s^2+s^3},\qquad
 \mathcal D=s\,\frac{d}{ds},\qquad
 \mathcal V=f(T^{30}).
\]
The marks \(P_\pm\) maintain channel orientation. Angular transport is the antecedent of the response; its functional reconstruction recovers this representation without replacing the memory archive from which it arises.

\begin{theorem}[Inversion of the constitutive response]
\label{cr28:thm:inversion-constitutiva} The function \(f\) is a decreasing bijection from \((0,1)\) onto \((3/4,1)\). Its inverse \(\psi\) assigns to \(r\) the unique root \(s\in(0,1)\) of
\begin{equation}
 r s^3+(r-1)(1+s+s^2)=0.
 \label{cr28:eq:cubic}
\end{equation}
The two labelled responses \(r_\pm\) determine
\begin{equation}
 s_\pm=\psi(r_\pm),\quad q_\pm=s_\pm^{1/30},\quad
 x=-\frac12\log(q_+q_-),\quad
 y=\frac12\log\frac{q_-}{q_+}.
 \label{cr28:eq:angular-recovery}
\end{equation}
For the complete operator \(S=\mathcal V^2\), positive with spectrum in \((9/16,1)\), functional calculus gives
\begin{equation}
 T=\bigl[\psi(S^{1/2})\bigr]^{1/30},\qquad
 \det T=\det\bigl[\psi(S^{1/2})\bigr]^{1/30}.
 \label{cr28:eq:operator-recovery}
\end{equation}
\end{theorem}
\begin{proof}
The negative derivative of \eqref{eq:vacuum-monotonic34} and the limits \(f(0^+)=1,\ f(1^-)=3/4\) prove bijectivity. Multiplying \(f(s)=r\) by its positive denominator gives \eqref{cr28:eq:cubic}. The positive root of order \(30\) recovers \(q_\pm\); the sum and difference of \(\log q_\pm=-x\mp y\) give the two angular coordinates. By positivity, \(S^{1/2}=\mathcal V\). Applying \(\psi\) to each spectral projector recovers \(T^{30}\) and the positive root gives \(T\). Taking its determinant completes the proof.
\end{proof}

This recovery of the complete operator does not extend to an isolated determinant. Indeed,
\[
 T_1=\operatorname{diag}(1/5,1/3),\qquad
 T_2=\operatorname{diag}(1/6,2/5)
\]
have determinant \(1/15\) but distinct spectra in the same contractive chamber. The proved injectivity implies \(f(T_1^{30})^2\ne f(T_2^{30})^2\). Both channels and their marks preserve recovery; reducing them to a product loses information.
% END CR28_CONSTITUTIVE_INVERSE

% BEGIN CR28_TRANSPORT_COMPOSITION
% Recuperación expositiva de INTERRELACIONES_ESTRUCTURALES_VACIO.md, §§6--8.
% Insertar después de la inversión cúbica y antes de las coordenadas logarítmicas.
\subsection{Transported composition of constitutive responses}
\label{cr28:comp-sec}

The angular transport and the incidences $90$ and $120$ originate in the
same APP--TRIT--TPK structure: the sheets retain evaluations and their
carries; the tritic regime retains orientation; the TPK calendar transports
phase and memory to the angular channels. The constitutive response
applies the function $f$ of \eqref{eq:vacuum-monotonic34} to those channels.
Its inversion expresses, in response coordinates, the composition of
transports sharing a common frame. The constructive order is thereby
preserved: transport, composition and spectral evaluation. The resulting
operation acts on this representation of the enriched state; its inverse
recovers the channels, while the complete memory history remains in the
underlying structure.

Let $s=q^{30}$ and $\psi=f^{-1}$. The exponent $30$ is the greatest
common divisor of the incidences $90$ and $120$. Cubic inversion determines
$\psi:(3/4,1)\longrightarrow(0,1)$. Adjoining the boundary response
$f(1)=3/4$ extends the bijection to
\[
 \psi:D\longrightarrow(0,1],\qquad D=[3/4,1).
\]
At this boundary the rational function $f$ is used; the quotient
$(1-q^{90})/(1-q^{120})$ has a removable singularity there.

\begin{proposition}[Response monoid and additive coordinate]
\label{cr28:comp-monoid}
The operation
\begin{equation}
 r\odot t=f\bigl(\psi(r)\psi(t)\bigr),\qquad r,t\in D,
 \label{cr28:comp-law}
\end{equation}
makes $D$ a commutative cancellative monoid with identity $3/4$.
The map
\begin{equation}
 \ell(r)=-\frac1{30}\log\psi(r)
 \label{cr28:comp-character}
\end{equation}
is an isomorphism from this monoid to $(\mathbb R_{\geq0},+)$.
The interval $(3/4,1)$ is closed under $\odot$; the identity belongs
to the boundary extension. The identity is the only invertible element
within $D$.
\end{proposition}
\begin{proof}
The product of two elements of $(0,1]$ belongs to the same interval, and
\[
 \psi(r\odot t)=\psi(r)\psi(t).
\]
For three responses, either association has image
$\psi(r)\psi(t)\psi(u)$; injectivity of $\psi$ proves associativity.
Commutativity follows from scalar multiplication. If
$r\odot t=r\odot u$, positivity of $\psi(r)$ permits cancellation,
giving $\psi(t)=\psi(u)$ and hence $t=u$. The parameter $1$ corresponds
to $3/4$ and provides the identity. The logarithm transforms products into
sums, and its range on $(0,1]$ shows that $\ell$ is a bijection onto
$\mathbb R_{\geq0}$. In particular, the response of a channel $q$
satisfies $\ell(r)=-\log q$. If both parameters are strictly less than
$1$, so is their product. Finally, two numbers in $(0,1]$ have product
$1$ only when both are $1$.
\end{proof}

\begin{proposition}[Admissible cancellation and inverse extension]
\label{cr28:comp-inverse}
For $r,u\in D$, the equation $r\odot t=u$ has a solution $t\in D$
exactly when $u\geq r$. This solution is unique and is given by
\begin{equation}
 t=f\!\left(\frac{\psi(u)}{\psi(r)}\right).
 \label{cr28:comp-cancellation}
\end{equation}
The same function $f$ is a bijection from $(0,\infty)$ onto $(0,1)$.
With this extension of $\psi$, the law~\eqref{cr28:comp-law} transports
the multiplicative group of positive parameters to the interval $(0,1)$.
Group inversion satisfies
\begin{equation}
 r^{\ominus}=\psi(r)r,\qquad
 r\odot r^{\ominus}=\frac34.
 \label{cr28:comp-group-inverse}
\end{equation}
\end{proposition}
\begin{proof}
The equation is equivalent to
$\psi(t)=\psi(u)/\psi(r)$. This quotient belongs to $(0,1]$ exactly
when $\psi(u)\leq\psi(r)$, which, by decreasing monotonicity, is
equivalent to $u\geq r$. The bijection proves existence and uniqueness.
The derivative in \eqref{eq:vacuum-monotonic34} is negative for every $s>0$;
the limits of $f$ at $0$ and infinity are $1$ and $0$, respectively.
This establishes the extended bijection. Multiplying the numerator and
denominator by $s^3$ gives
\[
 f(s^{-1})=\frac{s^3+s^2+s}{s^3+s^2+s+1}=s f(s).
\]
The multiplicative inverse of $s=\psi(r)$ therefore has response
$\psi(r)r$. Its composition with $r$ corresponds to the parameter $1$.
\end{proof}

The extension to $s>1$, equivalently to $q>1$, is the algebraic extension
of the reader. The chapter's physical contractive sector retains its
declared domain. The coordinate $\ell$ extends to an isomorphism with
$(\mathbb R,+)$ and changes sign under group inversion.

\subsubsection{Operator realization in a common sheet frame}

\begin{proposition}[Composition with marked projectors]
\label{cr28:comp-operator}
Let $\mathcal R=\mathcal R^*$, $\mathcal R^2=I$ and
$P_\pm=(I\pm\mathcal R)/2$. For $j=a,b$, let
\[
 T_j=q_{j,+}P_++q_{j,-}P_-,\qquad
 \mathcal V_j=f(T_j^{30}),\qquad 0<q_{j,\pm}<1.
\]
The composition $T_{a\circ b}=T_aT_b$ satisfies
\begin{equation}
 \mathcal V_{a\circ b}
 =f\bigl(\psi(\mathcal V_a)\psi(\mathcal V_b)\bigr).
 \label{cr28:comp-operator-law}
\end{equation}
Its two responses are $r_{a,\pm}\odot r_{b,\pm}$. If
$T_j=\exp[-(x_jI+y_j\mathcal R)]$ with $x_j>|y_j|$, then
\begin{equation}
 T_aT_b=
 \exp[-((x_a+x_b)I+(y_a+y_b)\mathcal R)],
 \label{cr28:comp-angular-addition}
\end{equation}
y $x_a+x_b>|y_a+y_b|$.
\end{proposition}
\begin{proof}
Orthogonality $P_+P_-=0$ gives
\[
 T_aT_b=q_{a,+}q_{b,+}P_++q_{a,-}q_{b,-}P_-.
\]
In particular, $T_a$ and $T_b$ commute. Functional calculus on these
projectors yields
$\psi(\mathcal V_j)=T_j^{30}$ and
$(T_aT_b)^{30}=T_a^{30}T_b^{30}$; applying $f$ proves
\eqref{cr28:comp-operator-law}. More explicitly, for every function $g$
defined on the two eigenvalues,
\[
 P_\pm g(T_j)=g(q_{j,\pm})P_\pm.
\]
Since $q_{j,\pm}=\exp[-(x_j\pm y_j)]$, their products prove
\eqref{cr28:comp-angular-addition}. The triangle inequality concludes
$x_a+x_b>|y_a|+|y_b|\geq|y_a+y_b|$.
\end{proof}

The recovered coordinates can be written directly as
\[
 x=\frac{\ell(r_+)+\ell(r_-)}2,\qquad
 y=\frac{\ell(r_+)-\ell(r_-)}2.
\]
Composition adds these pairs. Simultaneous sheet exchange permutes the
two channels of each transport, acts as $(x,y)\mapsto(x,-y)$ and is an
automorphism of the componentwise law. Inversion of both transports in
the algebraic extension acts as $(x,y)\mapsto(-x,-y)$. The two
involutions commute and retain distinct functions: one exchanges
orientation markings; the other inverts the transport.

To represent exchange by conjugation, use an operator $L$ satisfying
$L\mathcal R L^{-1}=-\mathcal R$. In the representation
\[
 \mathcal R=\begin{pmatrix}0&1\\1&0\end{pmatrix},\qquad
 L=\begin{pmatrix}1&0\\0&-1\end{pmatrix},
\]
direct multiplication gives $LP_\pm L^{-1}=P_\mp$, and hence
$LT_jL^{-1}=q_{j,-}P_++q_{j,+}P_-$. Conjugation by $\mathcal R$
preserves its own projectors. A trace combining the channels likewise
cannot replace the projectors in this functional identity: their labels
contribute to the composition.

The common-frame hypothesis is substantive. For example, the positive
definite operators
\[
 A_0=\begin{pmatrix}1/2&0\\0&1/3\end{pmatrix},\qquad
 B_0=\begin{pmatrix}1/2&1/6\\1/6&1/2\end{pmatrix}
\]
have positive eigenvalues, but
\[
 A_0B_0=\begin{pmatrix}1/4&1/12\\1/18&1/6\end{pmatrix}
\]
is not self-adjoint. The proposition specifies transport composition on
common projectors; its domain does not encompass an arbitrary mixture of
material media. Realizing a concatenation of HMT routes additionally
retains their memory data: algebraic closure of the parameter cone
describes the representation and does not by itself select new primary
routes.

\subsubsection{Channel product and transport composition}

The normalized speed of one state remains determined by
\[
 \widehat c=(r_+r_-)^{-1}.
\]
Here the ordinary product combines the two responses of that same state.
The law $\odot$ evaluates a different operation: composition of two
transports on each sheet. The following exact calculation distinguishes
their arguments:
\begin{equation}
 f(1/2)=\frac{14}{15},\qquad
 \frac{14}{15}\odot\frac{14}{15}=f(1/4)=\frac{84}{85}
 \ne\frac{196}{225}=\left(\frac{14}{15}\right)^2.
 \label{cr28:comp-rational-example}
\end{equation}
The fractions in this example are algebraic test parameters. The
expression for $\widehat c$ and the preceding constitutive identities
remain unchanged. The added content is an explicit law transporting
composition and recovering its additive angular coordinate.

% END CR28_TRANSPORT_COMPOSITION

% BEGIN CR28_CATALAN_RECOVERY
\subsection{Functional reconstruction of the oriented moment}
\label{cr28:sec:catalan-recovery}

The differential relation~\eqref{cr28:eq:catalan-problem} admits an inverse preserving the complete system of moments. Its domain is the constitutive function produced by the cycle sectors, with the orientation of~\eqref{eq:vacuum-angular-channels}. Incidences $90=3\cdot30$ and $120=4\cdot30$ fix that function before its evaluation in channels $s_\pm=q_\pm^{30}$. The condition at the origin corresponds to disappearance of the moment when its contractive argument vanishes. This condition also determines the integration constants.

\begin{theorem}[Integral reconstruction and uniqueness]
\label{cr28:thm:catalan-recovery}
Let $f(s)=(1-s^3)/(1-s^4)$, for $0<s<1$, be the response
of~\eqref{eq:three-four-completion}, and let $\mathcal D=s\,d/ds$.
There is a unique function $F\in C^2((0,1))$ satisfying
\begin{equation}
 \mathcal D^2F(s)=\frac{s(1+s)}{1+s+s^2}f(s),
 \qquad \lim_{s\downarrow0}F(s)=0.
 \label{cr28:eq:catalan-problem}
\end{equation}
It is given by
\begin{align}
 F(s)&=\int_0^s\log\frac{s}{t}\,
          \frac{1+t}{1+t+t^2}f(t)\,dt
       =\int_0^s\frac{\log(s/t)}{1+t^2}\,dt
       =\mathcal C_2(s),
 \label{cr28:eq:catalan-integral}\\
 f(s)&=\frac{1+s+s^2}{s(1+s)}\mathcal D^2F(s).
 \label{cr28:eq:catalan-inverse}
\end{align}
Its continuous extension to $s=1$ recovers $F(1)=G_{\rm Cat}$.
\end{theorem}
\begin{proof}
The factorization $1+t+t^2+t^3=(1+t)(1+t^2)$ gives
\[
 \frac{1+t}{1+t+t^2}f(t)=\frac1{1+t^2}.
\]
The integrand of~\eqref{cr28:eq:catalan-integral} is
integrable at zero, since
\[
 0\leq F(s)\leq\int_0^s\log(s/t)\,dt=s.
\]
Hence $F(s)\to0$. Differentiation on the open interval,
first on a truncated interval and then by dominated convergence,
gives
\[
 \mathcal DF(s)=\int_0^s\frac{dt}{1+t^2},
 \qquad \mathcal D^2F(s)=\frac{s}{1+s^2}
     =\frac{s(1+s)}{1+s+s^2}f(s).
\]
In particular, $F$ is of class $C^2$ and satisfies the equation.
If $F_1,F_2$ are two solutions, their difference $H=F_1-F_2$
satisfies $\mathcal D^2H=0$. Setting $z=\log s$ gives
$d^2H(e^z)/dz^2=0$, so $H(s)=a\log s+b$.
The existence of a zero limit as $s\downarrow0$ first forces
$a=0$ and then $b=0$. The condition on $F$ therefore determines
the solution and also the limit $\mathcal DF(s)\to0$.

For $s<1$, the geometric expansion of $(1+t^2)^{-1}$ and
absolute integrability of its weighted terms allow termwise integration:
\[
 \int_0^s t^{2k}\log(s/t)\,dt
     =\frac{s^{2k+1}}{(2k+1)^2},\qquad
 F(s)=\sum_{k\geq0}\frac{(-1)^ks^{2k+1}}{(2k+1)^2}.
\]
The series represents the integral system~\eqref{cr28:eq:catalan-integral}. Absolute and uniform convergence of this last series on $[0,1]$ permits passage to the endpoint $s=1$. In the integral, this is also permitted by the integrable bound $\log(1/t)$, extending the integrand by zero for $t>s$. Finally, solving the differential equation for $f$ gives \eqref{cr28:eq:catalan-inverse}, with positive denominators on $(0,1)$.
\end{proof}

The limit of the moment and the limit of its second derivative are
taken in their respective convergence classes. In particular,
\[
 \lim_{s\uparrow1}\mathcal D^2\mathcal C_2(s)=\frac12
\]
is the radial Abel limit of the differentiated series. Evaluation
of $\mathcal C_2(1)$ directly uses its absolutely convergent
series. This distinction preserves the same depth operator in
the interior and in its boundary readout.

\subsection{Compatibility of reconstruction with channels and speed}
\label{cr28:sec:channels-velocity}

The theorem restores a function from another previously constructed function. Cubic inversion~\eqref{cr28:eq:cubic} instead acts on evaluations $r_\pm$: it recovers arguments $s_\pm$ within that fixed system. Composition of the two operations preserves quarter-turn orientation, labels $P_\pm$ and transport incidences. Thus $G_{\rm Cat}=\mathcal C_2(1)$ is the endpoint reading of a functional object also participating in the interior responses $r_\pm=f(s_\pm)$.

For $\mathcal V_m=\mathcal V\otimes I_{9^m}$ and
$\tau_m=\operatorname{Tr}/(2\cdot9^m)$, the speed readout
retains exactly the normalization of
\eqref{eq:vacuum-traces}:
\begin{equation}
 \exp[-2\tau_m(\log\mathcal V_m)]
   =\frac1{r_+r_-}
   =(\det\mathcal V)^{-1}=\widehat c.
 \label{cr28:eq:velocity-normalized}
\end{equation}
Indeed, each eigenvalue $r_\pm$ is repeated $9^m$ times, so
$2\tau_m(\log\mathcal V_m)=\log r_++\log r_-$.
The identity uses the determinant of the two-sheet realization
and the normalized trace of its amplification. The previously fixed
dimensional sections then give
\[
 c_{\rm int}=c_{\rm vac}=\widehat c\,u_C,
 \qquad \ell_0=c_{\rm vac}t_0,
\]
with the same state, chart and elementary duration as~\eqref{cr28:eq:same-velocity}. Functional restoration preserves this velocity section when composing it with action and length; spectral multiplicity is absorbed by the declared normalization.

% END CR28_CATALAN_RECOVERY

% BEGIN CR28_CONSTITUTIVE_RADIAL_LINK
\subsection{A common constitutive section in action, length and gravity}
\label{cr28:sec:constitutive-radial} Relations \eqref{eq:vacuum-sections} and the preceding restoration maintain a unique velocity section:
\begin{equation}
 c_{\rm int}=c_{\rm vac}=\widehat c\,u_C,\qquad
 \ell_0=c_{\rm int}t_0,\qquad
 c_{\rm int}=(\mu_{\rm vac}\varepsilon_{\rm vac})^{-1/2}.
 \label{cr28:eq:same-velocity}
\end{equation}
The final equality is an identity of sections in compatible dimensional bases. The radial proof of \cref{g26:sec:rigidez-radial-es} first constructs \(L=(5759/23040)\alpha^{16}L_*\) and then determines
\begin{equation}
 G_{\rm ret}=\frac{c_{\rm int}^3L^2}{\hbar_{\rm ret}}
 =\frac{L^2}
 {\hbar_{\rm ret}(\mu_{\rm vac}\varepsilon_{\rm vac})^{3/2}}.
 \label{cr28:eq:constitutive-gravity}
\end{equation}
The second expression results from substituting the constitutive identity into the first; it introduces no new calibration of either channel. In coordinates, if \(c_{\rm int}=c_*\widehat c\), then \(G_{\rm ret}=c_*^3\widehat c^{\,3}L^2/\hbar_{\rm ret}\). The propagation factor enters only once. Catalan recovery, composed transport and radial response thus preserve the same angular origin, its marks and its normalization.
% END CR28_CONSTITUTIVE_RADIAL_LINK

